\documentclass[12pt, letterpaper]{article}
\usepackage[utf8]{inputenc}
\usepackage[spanish]{babel}
\usepackage{geometry}
\usepackage{enumitem}

\title{
    \Large \textbf{Facultad de Ciencias} \\[0.3em]
    \large \textsc{Manejo de Datos} \\[1.5em]
    \LARGE \textbf{Buenas Prácticas, POO y Principios SOLID en Python} \\[0.5em]
    \Large Tarea 1
}
\author{
    \textbf{Autores:} \\[0.5em]
    García Galindo Melissa \\
    Martínez Molina Karen \\
    Montoya García Ángel Eduardo \\
    Vázquez García Mauricio Alexis
}
\date{Fecha de entrega: 8 de Septiembre del 2026}

\begin{document}
\maketitle
\newpage

\section{Investigación}

\subsection{Parte A - PEP 8}
PEP 8 (\textit{Python Enhancement Proposal 8}) se refiere a una guía de convenciones estilísticas a seguir para escribir código en python, que van desde conocer el tamaño máximo para las lineas de código, propuestas en como nombrar a las variables, volver el código legible, entre otras cosas. Incluso existen \textbf{linters} (herramientas que analizan el código de forma automática para detectar errores o malas practicas) y \textbf{autoformatters} (herramientas que pueden reescribir código o indicar nuestros errores) que podemos usar en caso de no recordar todas las convenciones al escribir código. A continuación colocamos 10 ejemplos donde se puedan visualizar las aplicaciones de estas convenciones al código:

\begin{enumerate}
    \item \textbf{Indentación consistente de 4 espacios} \\
    \textbf{Principio:} Emplear exactamente 4 espacios por nivel de anidación. No utilizar tabuladores ni mezclar tabuladores con espacios.
    \begin{itemize}
        \item \textit{Incorrecto:}
\begin{verbatim}
def calcular_area(base, altura):
  area = (base * altura) / 2
        return area
\end{verbatim}
        \item \textit{Correcto:}
\begin{verbatim}
def calcular_area(base, altura):
    area = (base * altura) / 2
    return area
\end{verbatim}
    \end{itemize}

    \item \textbf{Longitud máxima de línea (79 caracteres)} \\
    \textbf{Principio:} Limitar todas las líneas de código a un máximo de 79 caracteres para evitar el desplazamiento horizontal y permitir la lectura simultánea de múltiples archivos.
    \begin{itemize}
        \item \textit{Incorrecto:}
\begin{verbatim}
def calcular_prima(suma_asegurada, factor_edad, descuento_no_fumador, recargo_extra):
    return (suma_asegurada * factor_edad / 1000) - descuento_no_fumador + recargo_extra
\end{verbatim}
        \item \textit{Correcto:}
\begin{verbatim}
def calcular_prima(suma_asegurada, factor_edad, 
                   descuento_no_fumador, recargo_extra):
    prima_base = (suma_asegurada * factor_edad) / 1000
    return prima_base - descuento_no_fumador + recargo_extra
\end{verbatim}
    \end{itemize}

    \item \textbf{Líneas en blanco para separación de estructuras} \\
    \textbf{Principio:} Separar funciones globales y clases de nivel superior con dos líneas en blanco. Separar métodos internos de una clase con una sola línea en blanco.
    \begin{itemize}
        \item \textit{Incorrecto:}
\begin{verbatim}
class Asegurado:
    def __init__(self, edad):
        self.edad = edad
    def es_mayor(self):
        return self.edad >= 18
def funcion_auxiliar():
    pass
\end{verbatim}
        \item \textit{Correcto:}
\begin{verbatim}
class Asegurado:
    def __init__(self, edad):
        self.edad = edad

    def es_mayor(self):
        return self.edad >= 18


def funcion_auxiliar():
    pass
\end{verbatim}
    \end{itemize}

    \item \textbf{Organización y orden de importaciones} \\
    \textbf{Principio:} Cada import debe residir en una línea independiente al inicio del archivo, agrupados en tres bloques: biblioteca estándar, librerías de terceros y módulos locales.
    \begin{itemize}
        \item \textit{Incorrecto:}
\begin{verbatim}
import sys, os, math
from aseguradora.modelos import Asegurado
\end{verbatim}
        \item \textit{Correcto:}
\begin{verbatim}
import os
import sys

from aseguradora.modelos import Asegurado
\end{verbatim}
    \end{itemize}

    \item \textbf{Espacios en blanco en expresiones y sentencias} \\
    \textbf{Principio:} Rodear operadores binarios (asignación, aritméticos, relacionales) con un espacio a cada lado, pero evitar espacios inmediatamente dentro de paréntesis, corchetes o llaves.
    \begin{itemize}
        \item \textit{Incorrecto:}
\begin{verbatim}
x=5
resultado = ( x * 2 ) + [ 1, 2, 3 ][ 0 ]
\end{verbatim}
        \item \textit{Correcto:}
\begin{verbatim}
x = 5
resultado = (x * 2) + [1, 2, 3][0]
\end{verbatim}
    \end{itemize}

    \item \textbf{Convenciones de nomenclatura (\textit{Naming Conventions})} \\
    \textbf{Principio:} Utilizar \texttt{snake\_case} para funciones y variables, \texttt{PascalCase} (o \textit{CapWords}) para nombres de clases, y \texttt{UPPER\_SNAKE\_CASE} para constantes globales.
    \begin{itemize}
        \item \textit{Incorrecto:}
\begin{verbatim}
tasa_cambio = 21.13
class factor_edad:
    def CalcularFactor(self, EdadAsegurado):
        pass
\end{verbatim}
        \item \textit{Correcto:}
\begin{verbatim}
TASA_CAMBIO = 21.13

class FactorEdad:
    def calcular_factor(self, edad_asegurado):
        pass
\end{verbatim}
    \end{itemize}

    \item \textbf{Comparaciones de identidad con \texttt{None}} \\
    \textbf{Principio:} Las comparaciones con valores únicos singleton como \texttt{None} deben efectuarse siempre empleando operadores de identidad (\texttt{is} o \texttt{is not}), jamás mediante igualdad de valor (\texttt{==}).
    \begin{itemize}
        \item \textit{Incorrecto:}
\begin{verbatim}
if valor == None:
    inicializar_variable()
\end{verbatim}
        \item \textit{Correcto:}
\begin{verbatim}
if valor is None:
    inicializar_variable()
\end{verbatim}
    \end{itemize}

    \item \textbf{Comprobaciones implícitas de verdad booleana} \\
    \textbf{Principio:} No comparar directamente secuencias vacías o valores booleanos contra \texttt{True} o \texttt{False}; evaluar la secuencia directamente por su valor de verdad intrínseco.
    \begin{itemize}
        \item \textit{Incorrecto:}
\begin{verbatim}
if len(lista_palabras) == 0:
    return None
if es_fumador == True:
    aplicar_recargo()
\end{verbatim}
        \item \textit{Correcto:}
\begin{verbatim}
if not lista_palabras:
    return None
if es_fumador:
    aplicar_recargo()
\end{verbatim}
    \end{itemize}

    \item \textbf{Manejo explícito de excepciones} \\
    \textbf{Principio:} Evitar el uso de cláusulas \texttt{except:} vacías que atrapen \texttt{BaseException} (incluyendo señales de terminación del sistema). Especificar la excepción concreta esperada.
    \begin{itemize}
        \item \textit{Incorrecto:}
\begin{verbatim}
try:
    edad = int(input("Edad: "))
except:
    print("Error al procesar.")
\end{verbatim}
        \item \textit{Correcto:}
\begin{verbatim}
try:
    edad = int(input("Edad: "))
except ValueError:
    print("Error: Ingrese un número entero válido.")
\end{verbatim}
    \end{itemize}

    \item \textbf{Espacios alrededor de argumentos por defecto} \\
    \textbf{Principio:} No utilizar espacios alrededor del signo \texttt{=} cuando se definan valores por defecto en argumentos de funciones o al pasar argumentos por palabra clave (\textit{kwargs}).
    \begin{itemize}
        \item \textit{Incorrecto:}
\begin{verbatim}
def normalizar_texto(texto, remover_acentos = True):
    return texto.strip()
\end{verbatim}
        \item \textit{Correcto:}
\begin{verbatim}
def normalizar_texto(texto, remover_acentos=True):
    return texto.strip()
\end{verbatim}
    \end{itemize}
\end{enumerate}

\newpage
\subsection{Parte B -- Principios de Diseño Orientado a Objetos (SOLID)}


\subsection{Referencias Consultadas}

\end{document}
